%% ma: L12-B06-0044
%% lop: 12
%% chuong: 2
%% bai: 6
%% dang: 12.6.19
%% loai: TN4
%% muc: NB
%% dap_an: "D"
%% nguon: "CĐ2.1.pdf giáo viên gửi 27/10/2026 (M07, Lê Minh Khôi – Nguyễn Anh Khôi, Chương II Chủ đề 1 Vectơ trong không gian), Câu 4"
%% kien_thuc: "Bài 6 (quy tắc ba điểm, quy tắc hiệu, quy tắc hình bình hành)"
%% trang_thai: cho-duyet
%% ly_do: "Dạng mới 12.6.19 đề xuất, chờ giáo viên duyệt"
\begin{ex}
Trong không gian, cho bốn điểm $A$, $B$, $C$, $D$ tạo thành hình bình hành $ABCD$. Mệnh đề nào sau đây sai?
\choice{$\overrightarrow{AC}=\overrightarrow{AB}+\overrightarrow{AD}$}{$\overrightarrow{AB}=\overrightarrow{DC}$}{$\overrightarrow{OA}+\overrightarrow{OC}=\vec 0$ với $O$ là tâm của hình bình hành}{\True $\overrightarrow{AC}=\overrightarrow{BD}$}
\loigiai{$\overrightarrow{AC}$ và $\overrightarrow{BD}$ là hai đường chéo, không cùng phương nên $\overrightarrow{AC}\ne\overrightarrow{BD}$. Các mệnh đề A, B, C đúng.}
\end{ex}
